% Options for packages loaded elsewhere
% Options for packages loaded elsewhere
\PassOptionsToPackage{unicode}{hyperref}
\PassOptionsToPackage{hyphens}{url}
\PassOptionsToPackage{dvipsnames,svgnames,x11names}{xcolor}
%
\documentclass[
  english,
  10pt,
  letterpaper,
  DIV=11,
  numbers=noendperiod]{scrartcl}
\usepackage{xcolor}
\usepackage[margin=1in]{geometry}
\usepackage{amsmath,amssymb}
\setcounter{secnumdepth}{5}
\usepackage{iftex}
\ifPDFTeX
  \usepackage[T1]{fontenc}
  \usepackage[utf8]{inputenc}
  \usepackage{textcomp} % provide euro and other symbols
\else % if luatex or xetex
  \usepackage{unicode-math} % this also loads fontspec
  \defaultfontfeatures{Scale=MatchLowercase}
  \defaultfontfeatures[\rmfamily]{Ligatures=TeX,Scale=1}
\fi
\usepackage{lmodern}
\ifPDFTeX\else
  % xetex/luatex font selection
\fi
% Use upquote if available, for straight quotes in verbatim environments
\IfFileExists{upquote.sty}{\usepackage{upquote}}{}
\IfFileExists{microtype.sty}{% use microtype if available
  \usepackage[]{microtype}
  \UseMicrotypeSet[protrusion]{basicmath} % disable protrusion for tt fonts
}{}
\makeatletter
\@ifundefined{KOMAClassName}{% if non-KOMA class
  \IfFileExists{parskip.sty}{%
    \usepackage{parskip}
  }{% else
    \setlength{\parindent}{0pt}
    \setlength{\parskip}{6pt plus 2pt minus 1pt}}
}{% if KOMA class
  \KOMAoptions{parskip=half}}
\makeatother
% Make \paragraph and \subparagraph free-standing
\makeatletter
\ifx\paragraph\undefined\else
  \let\oldparagraph\paragraph
  \renewcommand{\paragraph}{
    \@ifstar
      \xxxParagraphStar
      \xxxParagraphNoStar
  }
  \newcommand{\xxxParagraphStar}[1]{\oldparagraph*{#1}\mbox{}}
  \newcommand{\xxxParagraphNoStar}[1]{\oldparagraph{#1}\mbox{}}
\fi
\ifx\subparagraph\undefined\else
  \let\oldsubparagraph\subparagraph
  \renewcommand{\subparagraph}{
    \@ifstar
      \xxxSubParagraphStar
      \xxxSubParagraphNoStar
  }
  \newcommand{\xxxSubParagraphStar}[1]{\oldsubparagraph*{#1}\mbox{}}
  \newcommand{\xxxSubParagraphNoStar}[1]{\oldsubparagraph{#1}\mbox{}}
\fi
\makeatother


\usepackage{longtable,booktabs,array}
\newcounter{none} % for unnumbered tables
\usepackage{calc} % for calculating minipage widths
% Correct order of tables after \paragraph or \subparagraph
\usepackage{etoolbox}
\makeatletter
\patchcmd\longtable{\par}{\if@noskipsec\mbox{}\fi\par}{}{}
\makeatother
% Allow footnotes in longtable head/foot
\IfFileExists{footnotehyper.sty}{\usepackage{footnotehyper}}{\usepackage{footnote}}
\makesavenoteenv{longtable}
\usepackage{graphicx}
\makeatletter
\newsavebox\pandoc@box
\newcommand*\pandocbounded[1]{% scales image to fit in text height/width
  \sbox\pandoc@box{#1}%
  \Gscale@div\@tempa{\textheight}{\dimexpr\ht\pandoc@box+\dp\pandoc@box\relax}%
  \Gscale@div\@tempb{\linewidth}{\wd\pandoc@box}%
  \ifdim\@tempb\p@<\@tempa\p@\let\@tempa\@tempb\fi% select the smaller of both
  \ifdim\@tempa\p@<\p@\scalebox{\@tempa}{\usebox\pandoc@box}%
  \else\usebox{\pandoc@box}%
  \fi%
}
% Set default figure placement to htbp
\def\fps@figure{htbp}
\makeatother


% definitions for citeproc citations
\NewDocumentCommand\citeproctext{}{}
\NewDocumentCommand\citeproc{mm}{%
  \begingroup\def\citeproctext{#2}\cite{#1}\endgroup}
\makeatletter
 % allow citations to break across lines
 \let\@cite@ofmt\@firstofone
 % avoid brackets around text for \cite:
 \def\@biblabel#1{}
 \def\@cite#1#2{{#1\if@tempswa , #2\fi}}
\makeatother
\newlength{\cslhangindent}
\setlength{\cslhangindent}{1.5em}
\newlength{\csllabelwidth}
\setlength{\csllabelwidth}{3em}
\newenvironment{CSLReferences}[2] % #1 hanging-indent, #2 entry-spacing
 {\begin{list}{}{%
  \setlength{\itemindent}{0pt}
  \setlength{\leftmargin}{0pt}
  \setlength{\parsep}{0pt}
  % turn on hanging indent if param 1 is 1
  \ifodd #1
   \setlength{\leftmargin}{\cslhangindent}
   \setlength{\itemindent}{-1\cslhangindent}
  \fi
  % set entry spacing
  \setlength{\itemsep}{#2\baselineskip}}}
 {\end{list}}
\usepackage{calc}
\newcommand{\CSLBlock}[1]{\hfill\break\parbox[t]{\linewidth}{\strut\ignorespaces#1\strut}}
\newcommand{\CSLLeftMargin}[1]{\parbox[t]{\csllabelwidth}{\strut#1\strut}}
\newcommand{\CSLRightInline}[1]{\parbox[t]{\linewidth - \csllabelwidth}{\strut#1\strut}}
\newcommand{\CSLIndent}[1]{\hspace{\cslhangindent}#1}

\ifLuaTeX
\usepackage[bidi=basic,shorthands=off]{babel}
\else
\usepackage[bidi=default,shorthands=off]{babel}
\fi
\ifLuaTeX
  \usepackage{selnolig} % disable illegal ligatures
\fi


\setlength{\emergencystretch}{3em} % prevent overfull lines

\providecommand{\tightlist}{%
  \setlength{\itemsep}{0pt}\setlength{\parskip}{0pt}}



 


\KOMAoption{captions}{tableheading}
\makeatletter
\@ifpackageloaded{caption}{}{\usepackage{caption}}
\AtBeginDocument{%
\ifdefined\contentsname
  \renewcommand*\contentsname{Table of contents}
\else
  \newcommand\contentsname{Table of contents}
\fi
\ifdefined\listfigurename
  \renewcommand*\listfigurename{List of Figures}
\else
  \newcommand\listfigurename{List of Figures}
\fi
\ifdefined\listtablename
  \renewcommand*\listtablename{List of Tables}
\else
  \newcommand\listtablename{List of Tables}
\fi
\ifdefined\figurename
  \renewcommand*\figurename{Figure}
\else
  \newcommand\figurename{Figure}
\fi
\ifdefined\tablename
  \renewcommand*\tablename{Table}
\else
  \newcommand\tablename{Table}
\fi
}
\@ifpackageloaded{float}{}{\usepackage{float}}
\floatstyle{ruled}
\@ifundefined{c@chapter}{\newfloat{codelisting}{h}{lop}}{\newfloat{codelisting}{h}{lop}[chapter]}
\floatname{codelisting}{Listing}
\newcommand*\listoflistings{\listof{codelisting}{List of Listings}}
\makeatother
\makeatletter
\makeatother
\makeatletter
\@ifpackageloaded{caption}{}{\usepackage{caption}}
\@ifpackageloaded{subcaption}{}{\usepackage{subcaption}}
\makeatother
\usepackage{bookmark}
\IfFileExists{xurl.sty}{\usepackage{xurl}}{} % add URL line breaks if available
\urlstyle{same}
% fallback for those not using the hyperref driver hyperxmp:
\makeatletter
\@ifundefined{xmpquote}{\newcommand{\xmpquote}[1]{#1}}{}
\makeatother
\hypersetup{
  pdftitle={American Calls and Puts under Black--Scholes--Merton from One-Driver Complex Geometric Brownian Motion},
  pdfauthor={Complex--Itô Random-Walk Research Project},
  pdflang={en},
  pdfkeywords={\xmpquote{American
options}, \xmpquote{Black--Scholes--Merton}, \xmpquote{dividend
yield}, \xmpquote{early exercise premium}, \xmpquote{free
boundary}, \xmpquote{short-expiry asymptotics}, \xmpquote{collocation}},
  colorlinks=true,
  linkcolor={blue},
  filecolor={Maroon},
  citecolor={Blue},
  urlcolor={Blue},
  pdfcreator={LaTeX via pandoc}}


\title{American Calls and Puts under Black--Scholes--Merton from
One-Driver Complex Geometric Brownian Motion}
\usepackage{etoolbox}
\makeatletter
\providecommand{\subtitle}[1]{% add subtitle to \maketitle
  \apptocmd{\@title}{\par {\large #1 \par}}{}{}
}
\makeatother
\subtitle{Phase-barrier geometry, a regime-aware boundary family, and a
collocation pricing formula --- Revision 2 after mathematical review}
\author{Complex--Itô Random-Walk Research Project}
\date{September 3, 2026}
\begin{document}
\maketitle
\begin{abstract}
\textbf{Revision 2} of the Phase 15 paper, correcting every item of the
external mathematical review of 2026-09-02. We price American calls and
puts under Black--Scholes--Merton assumptions with continuous dividend
yield \(q\) as European Black--Scholes plus the early-exercise premium
(EEP; Kim, Jacka, Carr--Jarrow--Myneni) evaluated on a collocated
boundary family. Revisions: (i) the put EEP is stated with its correct
explicit sign (the one-line unified form needs an outer factor \(\phi\)
for puts); (ii) the short-expiry boundary asymptotics are the
established three regimes of Evans--Kuske--Keller and Chen--Zhu ---
which our solved boundaries reproduce quantitatively --- not a single
universal law; (iii) the boundary family is redesigned to be
dimensionless (declared logarithm scale \(\hat\tau\)), smooth on
\((0,\infty)\), regime-aware, and monotone through the former artificial
cusp at \(\tau=1\); (iv) all stationarity/envelope claims are removed
--- the substituted-boundary EEP functional is first-order sensitive to
boundary error, demonstrated by perturbation; (v) every benchmark is
recomputed against a converged Broadie--Detemple adjacent-averaged tree
oracle, with the Volterra level-set solver demoted to a method under
test and its grid sensitivity quantified; (vi) metric labels corrected
(RMSE / true MAE / maximum absolute error). On a 90-configuration grid
(450 prices, \(T\le3\), \(\sigma\le0.6\)) the revised formula attains
RMSE \(1.5\times10^{-3}\), mean absolute error \(8\times10^{-4}\),
maximum error \(7.8\times10^{-3}\) against the oracle, versus
\(0.279/0.162/1.131\) for Barone-Adesi--Whaley. No closed form for the
true free boundary is claimed.
\end{abstract}

\renewcommand*\contentsname{Table of contents}
{
\setcounter{tocdepth}{1}
\tableofcontents
}

\section{Revision record}\label{sec-revision}

This is Revision 2 of the Phase 15 manuscript, responding point by point
to \texttt{phase15-mathematical-review.md} (2026-09-02):

\begin{enumerate}
\def\labelenumi{\arabic{enumi}.}
\tightlist
\item
  \textbf{Put sign.} Section~\ref{sec-eep} states the call and put
  premiums separately and notes the outer factor \(\phi\) required by
  any unified one-line form.
\item
  \textbf{Short-expiry asymptotics.} Section~\ref{sec-asymptotics}
  replaces the universal claim with the three-regime results of Evans et
  al. (2001) and Chen and Zhu (2010), and shows our numerics reproduce
  them.
\item
  \textbf{Boundary family.} Section~\ref{sec-family} removes the
  \(|\ln\max(\tau,\varepsilon)|\) term (artificial cusp at \(\tau=1\),
  dimensionful logarithm, nonmonotone candidates) and replaces it with
  smooth dimensionless regime-aware bases.
\item
  \textbf{Stationarity.} All envelope/stationarity claims are deleted;
  Section~\ref{sec-collocation} demonstrates first-order sensitivity of
  the substituted-boundary functional.
\item
  \textbf{Reference solver.} Section~\ref{sec-solver} documents the
  solver's flat-step fallbacks and grid sensitivity, and
  Section~\ref{sec-results} benchmarks against a converged tree oracle
  instead.
\item
  \textbf{Metric labels.} RMSE / true MAE / MaxAE throughout; the worst
  configurations are identified explicitly.
\end{enumerate}

\section{Setup and exact machinery}\label{sec-setup}

\subsection{Risk-neutral dynamics and the complex-GBM
picture}\label{sec-geometry}

Under the risk-neutral measure, \(dS_t/S_t=(r-q)\,dt+\sigma\,dW_t\), and
the Phase 4 one-driver complex process (Complex--Itô Random-Walk
Research Project 2026)

\begin{equation}\protect\phantomsection\label{eq-complex-gbm}{
\mathcal Z_t
=
\mathcal Z_0
\exp\!\left(
\left[\left(r-q-\tfrac12\sigma^2\right)+i\omega\right]t
+(\sigma+i\beta)W_t
\right)
}\end{equation}

has modulus \(|\mathcal Z_t|=S_t\) (with \(|\mathcal Z_0|=S_0\)) and
unwrapped phase \(\Theta_t=\Theta_0+\omega t+\beta W_t\). In the
drift-matched gauge the phase is a deterministic affine function of
log-price, \(\Theta_t=\Theta_0+(\beta/\sigma)\ln(S_t/S_0)\), so the
exercise boundary becomes a moving phase barrier. \textbf{Scope, stated
plainly:} this is an invertible reparameterization (for
\(\beta/\sigma>0\); the exercise inequality reverses for \(\beta<0\),
and at \(\beta=0\) phase carries no spot information). It does not
reduce the optimal-stopping problem or add pricing information by
itself; the substance of this paper is the boundary family and the EEP
machinery that follow.

\subsection{The early-exercise premium, with correct put
sign}\label{sec-eep}

Let \(\tau=T-t\) and \(b(\tau)\) the boundary in time to maturity. With
\(d_{1,2}(S,B,s)=[\ln(S/B)+(r-q\pm\frac12\sigma^2)s]/(\sigma\sqrt s)\),
the call and put premiums are

\begin{equation}\protect\phantomsection\label{eq-eep-call}{
\mathrm{EEP}^{c}(S,\tau)
=
\int_0^\tau
\Big[
qS e^{-qs}\Phi\big(d_1(S,b(\tau-s),s)\big)
-rK e^{-rs}\Phi\big(d_2(S,b(\tau-s),s)\big)
\Big]ds,
}\end{equation}

\begin{equation}\protect\phantomsection\label{eq-eep-put}{
\mathrm{EEP}^{p}(S,\tau)
=
\int_0^\tau
\Big[
rK e^{-rs}\Phi\big(-d_2(S,b(\tau-s),s)\big)
-qS e^{-qs}\Phi\big(-d_1(S,b(\tau-s),s)\big)
\Big]ds.
}\end{equation}

Equivalently, the unified form carries an \textbf{outer factor \(\phi\)}
for puts:

\begin{equation}\protect\phantomsection\label{eq-eep-unified}{
V_A=V_E+\phi\int_0^\tau
\Big[
qS e^{-qs}\Phi(\phi d_1)-rK e^{-rs}\Phi(\phi d_2)
\Big]ds,\qquad \phi=\begin{cases}+1&\text{call}\\-1&\text{put.}\end{cases}
}\end{equation}

Omitting the factor gives a put \emph{below} its European price (base
case: \(7.493\) versus the correct \(7.661\) at \(r=q=0.05\),
\(\sigma=0.2\), \(T=1\)), an arbitrage violation; the implementation and
a unit test enforce the correct sign. The value-matching boundary
equation inherits the same factor.

Anchors (Kim 1990; Jacka 1991): \(b_c(0^+)=K\max(1,r/q)\) (infinite for
\(q=0\)), \(b_p(0^+)=K\min(1,r/q)\); perpetual boundaries
\(b_\infty^c=K s_+/(s_+-1)\), \(b_\infty^p=K s_-/(s_-1)\) from the roots
of \(\frac12\sigma^2s(s-1)+(r-q)s-r=0\) (\(s_+>1\) when \(q>0\); at
\(q=0\), \(s_+=1\) and the call boundary is infinite, i.e.~never
exercise, recovering Merton (1973)).

\section{Short-expiry boundary asymptotics: three
regimes}\label{sec-asymptotics}

The established results for the put boundary near expiry (Evans et al.
2001; Chen and Zhu 2010, eqs. (3.15)--(3.22)) are

\begin{equation}\protect\phantomsection\label{eq-regimes}{
\begin{aligned}
q<r:&\quad K-b_p(\tau)\sim K\sigma\sqrt{\tau|\ln\tau|},\\
q=r:&\quad K-b_p(\tau)\sim K\sigma\sqrt{2\tau|\ln\tau|},\\
q>r:&\quad b_p(\tau)=\tfrac rq K\left[1-\xi_1\sigma\sqrt{2\tau}+O(\tau)\right],
\qquad \xi_1\approx0.4517,
\end{aligned}
}\end{equation}

with calls inheriting mirrored regimes under \(r\leftrightarrow q\).
Fitting our solved near-expiry boundaries (flat fallback steps excluded
and counted):

\begin{longtable}[]{@{}
  >{\raggedright\arraybackslash}p{(\linewidth - 8\tabcolsep) * \real{0.0808}}
  >{\raggedright\arraybackslash}p{(\linewidth - 8\tabcolsep) * \real{0.4545}}
  >{\raggedright\arraybackslash}p{(\linewidth - 8\tabcolsep) * \real{0.1212}}
  >{\raggedright\arraybackslash}p{(\linewidth - 8\tabcolsep) * \real{0.2222}}
  >{\raggedright\arraybackslash}p{(\linewidth - 8\tabcolsep) * \real{0.1212}}@{}}
\caption{Near-expiry regime fits from the Volterra
solutions.}\label{tbl-regimes}\tabularnewline
\toprule\noalign{}
\begin{minipage}[b]{\linewidth}\raggedright
Regime
\end{minipage} & \begin{minipage}[b]{\linewidth}\raggedright
Fit (\(\sqrt{2\tau|\ln\tau|}\) normalization)
\end{minipage} & \begin{minipage}[b]{\linewidth}\raggedright
Prediction
\end{minipage} & \begin{minipage}[b]{\linewidth}\raggedright
\(\sqrt{2\tau}\) slope
\end{minipage} & \begin{minipage}[b]{\linewidth}\raggedright
Prediction
\end{minipage} \\
\midrule\noalign{}
\endfirsthead
\toprule\noalign{}
\begin{minipage}[b]{\linewidth}\raggedright
Regime
\end{minipage} & \begin{minipage}[b]{\linewidth}\raggedright
Fit (\(\sqrt{2\tau|\ln\tau|}\) normalization)
\end{minipage} & \begin{minipage}[b]{\linewidth}\raggedright
Prediction
\end{minipage} & \begin{minipage}[b]{\linewidth}\raggedright
\(\sqrt{2\tau}\) slope
\end{minipage} & \begin{minipage}[b]{\linewidth}\raggedright
Prediction
\end{minipage} \\
\midrule\noalign{}
\endhead
\bottomrule\noalign{}
\endlastfoot
put, \(r=0.07,q=0.03\) (\(q<r\)) & 0.724 & \(1/\sqrt2=0.707\) & --- &
--- \\
put, \(r=q=0.05\) & 1.044 & \(1\) & --- & --- \\
put, \(r=0.03,q=0.07\) (\(q>r\)) & 0.09 (window artifact) & fits
\(\sqrt\tau\), not \(\sqrt{\tau|\ln\tau|}\) & 0.0393 &
\((r/q)\xi_1\sigma=0.0387\) \\
\end{longtable}

The earlier Phase 15 interpretation (``universal shape with
\(r-q\)-dependent coefficient'') is withdrawn: the measured values are
precisely the known regimes, including the fact that fitting a
\(\sqrt{\tau|\ln\tau|}\) regressor to the \(q>r\) square-root law
produces an apparent coefficient that tends to zero as
\(\tau\downarrow0\).

\section{The boundary family (redesigned)}\label{sec-family}

With \(s_\phi=+1\) for calls and \(-1\) for puts, write the signed
coefficients \(\eta_{\log}=s_\phi h_{\log}\),
\(\eta_{\mathrm{root}}=s_\phi
h_{\mathrm{root}}\) with \(h_{\log},h_{\mathrm{root}}\ge0\):

\begin{equation}\protect\phantomsection\label{eq-family}{
\boxed{
\ln b(\tau)=\ln b_\infty+\ln\!\left(\frac{b_0}{b_\infty}\right)e^{-\kappa_1\tau}
+\sigma e^{-\kappa_2\tau}\Big[
\eta_{\log}\,\psi_{\log}(\tau;\hat\tau)
+\eta_{\mathrm{root}}\,\psi_{\mathrm{root}}(\tau)
\Big]
}
}\end{equation}

\begin{equation}\protect\phantomsection\label{eq-shapes}{
\psi_{\log}(\tau;\hat\tau)=\sqrt{2\tau\ln\!\left(1+\frac{\hat\tau}{\tau}\right)},
\qquad
\psi_{\mathrm{root}}(\tau)=\sqrt{2\tau}.
}\end{equation}

Properties, each enforced by a unit test:

\begin{itemize}
\tightlist
\item
  \textbf{Dimensionless.} \(\hat\tau\) is a declared time scale in the
  same unit as \(\tau\); changing units rescales both, and the price is
  unit-invariant (Section~\ref{sec-identities}).
\item
  \textbf{Smooth.} \(\psi_{\log}\) is smooth on \((0,\infty)\);
  \(\psi_{\log}\sim\sqrt{2\tau|\ln(\tau/\hat\tau)|}\) as
  \(\tau\downarrow0\). The former term \(|\ln\max(\tau,\varepsilon)|\)
  had an infinite-slope cusp at \(\tau=1\) --- an arbitrary,
  dimension-dependent point --- which is gone.
\item
  \textbf{Regime-aware.} The two bases carry the logarithmic and
  square-root short-expiry shapes of Eq.~\ref{eq-regimes}; collocation
  allocates their weights.
\item
  \textbf{Anchored.} \(b(0^+)=b_0\) and \(b(\tau)\to b_\infty\) for all
  parameter values.
\item
  \textbf{Monotone in practice.} The old cusp produced
  \(b(0.9)=132.3>b(1.0)=129.9<b(1.1)=136.2\) for the \(T=3\),
  \(r=q=0.05\), \(\sigma=0.2\) call; the redesigned family gives a
  monotone boundary through that neighborhood (regression-tested).
  Monotonicity is a numerically verified property of the collocated
  members, not an analytic guarantee of the whole family.
\end{itemize}

\section{Collocation: a reported projection, no
stationarity}\label{sec-collocation}

\((\kappa_1,\kappa_2,h_{\log},h_{\mathrm{root}})\) \textbf{minimize} the
sum of squared value-matching residuals at four fractional nodes
(\(\tau\in\{0.2,0.4,0.6,0.8\}\cdot T\)) under box constraints,
multi-start L-BFGS-B. This is a projection: residuals are generally
\emph{not} zero, and the published record reports, per calibration, the
objective, the four residuals, the optimizer status, and the multi-start
sensitivity. For the base case the residuals are
\(\sim10^{-5}\)--\(10^{-4}\); degenerate long-maturity cases can leave
residuals at the \(10^{-2}\) level, which is reported, not hidden.

\textbf{No stationarity claim.} The substituted-boundary functional
\(V_E+\mathrm{EEP}[b]\) is first-order sensitive to boundary error: for
\(d_{1,2}\) as above, the identity
\(Se^{-qs}\varphi(d_1)=Be^{-rs}\varphi(d_2)\) gives
\(\partial\,\mathrm{integrand}/\partial B
=e^{-rs}\varphi(d_2)(rK-qB)/(B\sigma\sqrt s)\), so the first variation
under a perturbation \(h\) is generally nonzero. Numerically, with
\(b_\varepsilon(u)=b(u)e^{\varepsilon\sin(\pi u/T)}\) on the base call,
\((V[b_\varepsilon]-V[b])/\varepsilon\) converges to \(\approx-0.345\)
as \(\varepsilon\to0\) through
\(10^{-2},5\times10^{-3},2.5\times10^{-3},
1.25\times10^{-3}\) --- first order, not second. Accuracy below is
established by direct benchmarking against the oracle, not by any
envelope argument.

\section{The Volterra solver and its known limits}\label{sec-solver}

The causal sequential level-set solver (Phase 15, revised root logic)
locates, at each grid point, where the value-matching residual reaches
the level \(-\delta K\) on the continuation side --- the exact residual
is identically zero on the exercise side (\(F\equiv0\) for \(y\ge b\)
for calls, \(y\le b\) for puts; the previously claimed put ``positive
hump'' is a numerical artifact and is no longer used). Its limits,
stated plainly and instrumented (\texttt{SOLVER\_DIAGNOSTICS}):

\begin{itemize}
\tightlist
\item
  \textbf{Flat fallbacks.} Points without a bracketed level crossing
  copy the previous value; counts are reported. Near \(\tau=0\) the
  equation loses traction (\(\partial F/\partial y\to0\)), and in
  long-maturity hard regimes the fallback fraction can exceed one half
  of all steps.
\item
  \textbf{Monotonicity is enforced} by clipping toward the perpetual
  limit: monotone output is a solver property, not an independent check.
\item
  \textbf{Grid sensitivity.} In the benchmark's hard regimes the solver
  is not a reliable reference (Table~\ref{tbl-grid}):
\end{itemize}

\begin{longtable}[]{@{}
  >{\raggedright\arraybackslash}p{(\linewidth - 2\tabcolsep) * \real{0.4286}}
  >{\raggedleft\arraybackslash}p{(\linewidth - 2\tabcolsep) * \real{0.5714}}@{}}
\caption{Grid behavior of the Volterra solver against the tree oracle.
The \(n=300\) value is an isolated outlier of \(\approx0.30\); fallback
counts exceed 200 steps in this
configuration.}\label{tbl-grid}\tabularnewline
\toprule\noalign{}
\begin{minipage}[b]{\linewidth}\raggedright
Settings
\end{minipage} & \begin{minipage}[b]{\linewidth}\raggedleft
Price (\(S=120\), \(r=0.05\), \(q=0.10\), \(\sigma=0.2\), \(T=3\), call)
\end{minipage} \\
\midrule\noalign{}
\endfirsthead
\toprule\noalign{}
\begin{minipage}[b]{\linewidth}\raggedright
Settings
\end{minipage} & \begin{minipage}[b]{\linewidth}\raggedleft
Price (\(S=120\), \(r=0.05\), \(q=0.10\), \(\sigma=0.2\), \(T=3\), call)
\end{minipage} \\
\midrule\noalign{}
\endhead
\bottomrule\noalign{}
\endlastfoot
Solver \(n=250\) & 20.634 \\
Solver \(n=300\) & 20.392 \\
Solver \(n=350\) & 20.677 \\
Solver \(n=700\) & 20.671 \\
Adjacent-averaged CRR \(N\in\{8000,10000,12000\}\) (mean) & 20.692
(spread \(1.1\times10^{-3}\)) \\
\end{longtable}

Consequently \textbf{all benchmarks use the converged Broadie--Detemple
adjacent-averaged tree} (Broadie and Detemple 1996) as the oracle
(\(N\in\{8000,12000\}\), mean; spread reported per row), with the solver
demoted to one method under test. (Prices also remain sensitive at the
\(10^{-2}\) level to environment-level floating-point differences, e.g.
numpy SIMD paths --- one more reason the tree, not the solver, is the
oracle.)

\section{Numerical results}\label{sec-results}

\subsection{Benchmarks against the
oracle}\label{benchmarks-against-the-oracle}

Reduced grid: 18 configurations \(\times\) 3 spots
(\(S\in\{90,100,110\}\)). Extended grid: 90 configurations \(\times\) 5
spots (\(S\in\{80,\ldots,120\}\), \(\sigma\in\{0.2,0.4,0.6\}\),
\(T\in\{0.5,1,3\}\), five \((r,q)\) pairs, both kinds). Metrics against
the oracle, labeled correctly:

\begin{longtable}[]{@{}lllll@{}}
\caption{Benchmark metrics against the converged tree
oracle.}\label{tbl-bench}\tabularnewline
\toprule\noalign{}
Grid & Method & RMSE & MAE (mean) & MaxAE \\
\midrule\noalign{}
\endfirsthead
\toprule\noalign{}
Grid & Method & RMSE & MAE (mean) & MaxAE \\
\midrule\noalign{}
\endhead
\bottomrule\noalign{}
\endlastfoot
Reduced & spiral formula & \textbf{0.00066} & \textbf{0.00043} &
\textbf{0.0028} \\
Reduced & Barone-Adesi--Whaley & 0.161 & 0.096 & 0.427 \\
Reduced & Volterra solver (\(n=300\)) & 0.0077 & 0.0039 & 0.031 \\
Extended & spiral formula & \textbf{0.00147} & \textbf{0.00081} &
\textbf{0.0078} \\
Extended & Barone-Adesi--Whaley & 0.279 & 0.162 & 1.131 \\
Extended & Volterra solver (\(n=300\)) & 0.0206 & 0.0070 & 0.301 \\
\end{longtable}

The worst spiral case on the extended grid is
\((r,q,\sigma,T)=(0.05,0.10,0.2,1)\) call at \(S=120\) (error
\(-0.0078\)); the worst BAW errors concentrate in deep-ITM long-maturity
puts, e.g. the \(T=3\), \(\sigma=0.2\), \(r=0.07\), \(q=0.03\) put at
\(S=80\) identified by the review. Raw formula values are used
throughout --- no clipping --- and deep-ITM raw values that sit slightly
outside no-arbitrage bounds are visible in the CSVs as diagnostic
information.

\subsection{Structural identities}\label{sec-identities}

\begin{itemize}
\tightlist
\item
  \textbf{Time-unit invariance.} The same option priced in years and in
  months (rates per month, \(\sigma/\sqrt{12}\), \(\hat\tau=12\) months)
  agrees to \(4.5\times10^{-5}\).
\item
  \textbf{Off-ATM put--call duality.} \(P(90,100;0.07,0.03)\) versus the
  swapped call \(C(100,90;0.03,0.07)\), each on its own solved boundary:
  agreement to \(2.0\times10^{-3}\).
\item
  \textbf{Zero-dividend call.} Returns the European price to machine
  precision.
\end{itemize}

\section{Scope and limitations}\label{sec-scope}

\begin{itemize}
\tightlist
\item
  \textbf{No closed form for the true boundary is claimed.} The
  constant-coefficient free-boundary problem remains open; this paper
  delivers a collocation formula on an engineered family, validated
  against a converged oracle. The honest one-line description: a
  benchmark-validated collocation formula built on standard EEP
  machinery and an invertible complex-phase reparameterization of
  log-price.
\item
  \textbf{Solver sensitivity.} The Volterra solver is grid-sensitive in
  hard regimes (Table~\ref{tbl-grid}) and carries a heavy flat-step
  fraction near maturity; it is not the oracle anywhere in this
  revision.
\item
  \textbf{Raw formula values.} No clipping: approximate-boundary prices
  can sit slightly outside no-arbitrage bounds; the deviations are
  reported, and the upper bounds \(C_A\le S\), \(P_A\le K\) are not
  enforced by the formula.
\item
  \textbf{Near-expiry fits.} The regime fits exclude flat fallback steps
  and are only as reliable as the remaining point count (reported
  alongside).
\item
  \textbf{Model scope.} Constant \(r,q,\sigma\), continuous dividends,
  standard single-asset American calls/puts.
\end{itemize}

\section{Reproducibility}\label{sec-repro}

All computations run in the existing \texttt{phase3-paper} conda
environment (Python 3.12, numpy, scipy, matplotlib); no packages
installed. Fixed seed \texttt{20260902}. Artifacts:

\begin{itemize}
\tightlist
\item
  \texttt{phase15\_american.py} --- EEP machinery, regime-aware spiral
  family, collocation, solver with diagnostics, BAW, tree oracle;
\item
  \texttt{tests/test\_phase15\_american.py} --- 22 tests including the
  review's requested coverage (put sign, off-ATM duality, regime
  asymptotics, time-unit invariance, monotonicity through \(\tau=1\), no
  hidden clipping, grid behavior);
\item
  \texttt{notebooks/build\_phase15\_american\_spiral.py} → executed
  notebook --- the computational record behind every number above
  (\texttt{phase15\_figures/}, \texttt{phase15\_tables/}, including the
  extended oracle CSV);
\item
  \texttt{paper/phase15/} --- this manuscript.
\end{itemize}

\section*{References}\label{references}
\addcontentsline{toc}{section}{References}

\protect\phantomsection\label{refs}
\begin{CSLReferences}{1}{1}
\bibitem[\citeproctext]{ref-broadiedetemple1996}
Broadie, Mark, and Jérôme Detemple. 1996. {``{A}merican Option
Valuation: New Bounds, Approximations, and a Comparison of Existing
Methods.''} \emph{Review of Financial Studies} 9 (4): 1211--50.

\bibitem[\citeproctext]{ref-chenzhu2010}
Chen, Wen-Ting, and Song-Ping Zhu. 2010. {``Optimal Exercise Price of
{A}merican Options Near Expiry.''} \emph{ANZIAM Journal} 51: 375--91.
\url{https://doi.org/10.1017/S1446181110000052}.

\bibitem[\citeproctext]{ref-phase4paper}
Complex--Itô Random-Walk Research Project. 2026. \emph{Native Complex
Geometric Brownian Motion: Itô Correction, Stochastic Rotation, and
Noise-Dimension Geometry}.

\bibitem[\citeproctext]{ref-evanskuskekeller2001}
Evans, Jonathan D., Randy Kuske, and Joseph B. Keller. 2001.
{``{A}merican Options on Assets with Dividends Near Expiry.''}
\emph{Mathematical Finance} 11 (2): 219--37.
\url{https://doi.org/10.1111/1467-9965.02008}.

\bibitem[\citeproctext]{ref-jacka1991}
Jacka, Saul D. 1991. {``Optimal Stopping and the {A}merican Put
Price.''} \emph{Annals of Applied Probability} 1 (1): 1--21.

\bibitem[\citeproctext]{ref-kim1990}
Kim, In Joon. 1990. {``The Analytic Valuation of American Options.''}
\emph{Review of Financial Studies} 3 (4): 547--72.

\bibitem[\citeproctext]{ref-merton1973}
Merton, Robert C. 1973. {``Theory of Rational Option Pricing.''}
\emph{Bell Journal of Economics and Management Science} 4 (1): 141--83.

\end{CSLReferences}




\end{document}
